% Source: Independent oriented tetrahedral sum; KACTL PolyhedronVolume.h cross-check.
\subsection{四面体与封闭多面体体积}
\label{geo:volume}
\AlgorithmInfo
    {求四面体体积及由定向三角面给出的多面体有向体积。}
    {四面体 $O(1)$；多面体 $O(F)$ 时间、$O(1)$ 额外空间。}
    {依赖三维点；表面须闭合、无自交，所有面法向一致朝外。}
\subsubsection*{原理}

四面体体积为 $|(b-a)\cdot((c-a)\times(d-a))|/6$。闭多面体可与任意参考点 o 构成一批有向四面体，累加
\[
 V=\frac16\sum_{(a,b,c)\text{ 为面}}(a-o)\cdot((b-o)\times(c-o)).
\]
内部公共部分按有向体积抵消。实现选择一个已有顶点作参考，减少大坐标误差，并用补偿求和。凸性不是必需条件，但每个面必须按从外侧看逆时针的方向列出；非三角面应先按其边界方向正确三角剖分。

\subsubsection*{使用教程}
\InlineCode{tetrahedron_volume(a,b,c,d)} 直接接收四个三维点，返回非负四面体体积；四点顺序不影响绝对体积，共面时理论值为 0，浮点结果须结合尺度判断。

\InlineCode{polyhedron_signed_volume(faces)} 接收三角面数组，每面为 \InlineCode{array<Point3,3>}，要求整个表面\textbf{闭合且定向一致}。从实体外侧看，每面的顶点应逆时针，使叉积法向量朝外；全部朝外得到正体积，全部反向得到负体积。确认闭合与方向后，取绝对值才是实体体积；若部分面反了，取绝对值也修不好。

函数不自动建三维凸包或校正面朝向，空面集返回 0；合法的凹多面体也可按闭合表面求和。下面用四个面描述同一个四面体，同时演示直接四点法；两函数均不改变输入，无预处理。
\begin{lstlisting}[style=mtycpp,numbers=none]
Point3 a{0,0,0},b{1,0,0},c{0,1,0},d{0,0,1};
ld v=tetrahedron_volume(a,b,c,d);
vector<array<Point3,3>> faces={{a,c,b},{a,b,d},{a,d,c},{b,c,d}};
ld closed=polyhedron_signed_volume(faces);
\end{lstlisting}
\needspace{7\baselineskip}
\subsubsection*{模板代码}
\CodeFile{code/geometry/volume.cpp}
